\DocumentMetadata{
  pdfversion=2.0,
  pdfstandard=ua-2,
  lang=en-US,
  tagging=on,
  tagging-setup={math/setup=mathml-SE,tabsorder=structure}
}
\documentclass[11pt,letterpaper]{article}
\usepackage[margin=0.68in,includefoot]{geometry}
\usepackage{newcomputermodern}
\usepackage{amsmath,amscd,mathtools}
\usepackage{tikz,adjustbox}
\providecommand{\square}{\mdlgwhtsquare}
\usepackage[hidelinks]{hyperref}
\hypersetup{
  pdftitle={Partial Differential Equations PhD Exam, September 2010},
  pdfauthor={University of Florida Department of Mathematics},
  pdfsubject={Graduate mathematics examination},
  pdfdisplaydoctitle=true,
  bookmarksopen=true
}
\setcounter{secnumdepth}{0}
\setlength{\parindent}{0pt}
\setlength{\parskip}{0.28em}
\setlength{\emergencystretch}{3em}
\setlength{\leftmargini}{2.15em}
\setlength{\leftmarginii}{2.65em}
\setlength{\leftmarginiii}{3em}
\renewcommand{\labelenumi}{\textbf{\theenumi.}}
\pagestyle{empty}
\raggedbottom
\allowdisplaybreaks
\definecolor{ProblemConcern}{HTML}{7A1F1F}
\newenvironment{problemconcern}{%
  \par\tagstructbegin{tag=Aside}\begingroup\color{ProblemConcern}
}{%
  \par\endgroup\tagstructend\nopagebreak[4]\vspace{0.12em}
}
\newenvironment{examproblems}[1]{%
  \begin{enumerate}%
  \setcounter{enumi}{#1}%
  \setlength{\topsep}{0.35em}%
  \setlength{\partopsep}{0pt}%
  \setlength{\itemsep}{0.48em}%
  \setlength{\parsep}{0.18em}%
}{\end{enumerate}}
\newenvironment{examsubparts}{%
  \begin{enumerate}%
  \setlength{\topsep}{0.18em}%
  \setlength{\partopsep}{0pt}%
  \setlength{\itemsep}{0.16em}%
  \setlength{\parsep}{0.1em}%
}{\end{enumerate}}
\newenvironment{examinstructions}{%
  \par\begingroup\itshape
}{%
  \par\endgroup\vspace{0.18em}
}
\makeatletter
\renewcommand\section{\@startsection{section}{1}{0pt}%
  {-1.25ex plus -.25ex minus -.1ex}%
  {0.55ex plus .1ex}%
  {\normalfont\large\bfseries}}
\renewcommand{\@maketitle}{%
  \newpage\null\vskip -1em%
  {\footnotesize\bfseries
    \href{https://www.ufl.edu/}{University of Florida}, \href{https://math.ufl.edu/}{Department of Mathematics}\par}%
  \vskip -0.5em%
  \tagstructbegin{tag=Title}%
    {\LARGE\bfseries\href{https://gma.math.ufl.edu/exams/pde-phd/}{Partial Differential Equations PhD Exam}, September 2010\par}%
  \tagstructend%
  \vskip 0.25em%
  \tagmcbegin{artifact}%
    \hrule height 1pt%
  \tagmcend%
  \vskip 1em%
}
\makeatother
\title{Partial Differential Equations PhD Exam, September 2010}
\author{}
\date{}
\begin{document}
\maketitle
\thispagestyle{empty}
\begin{examproblems}{0}
\item
Let~\(U\) be either \(\mathbb R_+^n\) or the unit ball in \(\mathbb R^n\) (\(n\ge 2\)). Write the expression of the Green’s function for \(-\Delta\) with respect to the Dirichlet condition on \(\partial U\).
\item
Let \(U\subset\mathbb R^n\) be a bounded open set with smooth boundary \(\partial U\), \(A=(a_{ij})_{n\times n}\) be a symmetric, positive definite constant matrix. For the following variational form: 
\[
I(w):=\frac12\int_U\sum_{i,j=1}^n a_{ij}\partial_iw\partial_jw\,dx;\qquad w\in H_g^1:=\{w\in H^1(U);\ w=g\text{ on }\partial U\},
\]
 where~\(g\) is a smooth function on \(\partial U\), prove that the minimizer~\(u\) exists and is smooth.
\item
Let \(A=(a_{ij})_{n\times n}\) be a symmetric and positive definite constant matrix. Suppose~\(u\) satisfies 
\[
\sum_{i,j=1}^n a_{ij}\partial_{ij}u=0,\qquad \text{in }\mathbb R^n,
\]
 and \(u\ge -1\) in \(\mathbb R^n\). Show that \(u\equiv\text{constant}\).
\item
Let \(\phi(x)=\dfrac{1}{n(n-2)\alpha(n)}|x|^{2-n}\) be the fundamental solution for the Laplace operator in \(\mathbb R^n\) (\(n\ge 3\)). Here \(\alpha(n)\) is the volume of the unit ball \(B_1\). Suppose \(f(x)\in C^2(\mathbb R^n)\) and satisfies 
\[
|D^jf(x)|\le C(1+|x|)^{-2-\epsilon-j},\qquad \forall x\in\mathbb R^n,\quad j=0,1,2,
\]
 where~\(\epsilon\) is a positive number. Then 
\[
u(x)=\int_{\mathbb R^n}\phi(x-y)f(y)\,dy
\]
 satisfies 
\[
-\Delta u(x)=f(x)\qquad \mathbb R^n.
\]
\item
Let \(u\in W^{1,p}(\mathbb R^n)\), where \(p>n\). Show that 
\[
\frac{1}{r^n}\int_{B(x,r)}|u(x)-u(y)|\,dy\le C(n)\int_{B(x,r)}\frac{|Du(y)|}{|x-y|^{n-1}}\,dy.
\]
\item
Write down an explicit formula for a solution of 
\[
\begin{cases}
u_t-\Delta u+cu=f, & \mathbb R^n\times(0,\infty),\\
u=g, & \mathbb R^n\times\{t=0\},
\end{cases}
\]
 where \(c\in\mathbb R\).
\item
\begin{problemconcern}
\textbf{Warning: Suspected error.} The source is internally inconsistent: it states “where \(g,h\) are smooth functions,” but~\(g\) does not occur in the displayed problem, while~\(f\) does. The stated data may also omit an initial-displacement condition. The intended correction is not uniquely determined, so the source reading has been retained.
\end{problemconcern}
Let \(U\subset\mathbb R^n\) be an open, bounded subset of \(\mathbb R^n\) with smooth boundary. Set \(U_T=U\times[0,T]\), \(\Gamma_T=\overline{U_T}\setminus U_T\), where \(T>0\). Prove that there exists at most one solution \(u\in C^2(\overline{U_T})\) of 
\[
\begin{cases}
u_{tt}-\Delta u=f, & U_T,\\
u=0, & \Gamma_T,\\
u_t=h, & U\times\{t=0\},
\end{cases}
\]
 where \(g,h\) are smooth functions.
\end{examproblems}
\end{document}
